\documentclass[11pt]{article}
\usepackage[utf8]{inputenc}
\usepackage[T1]{fontenc}
\usepackage{lmodern}
\usepackage[margin=1in]{geometry}
\usepackage{amsmath,amssymb}
\usepackage{textcomp}
\usepackage[dvipsnames]{xcolor}
\usepackage{enumitem}
\usepackage{parskip}
\usepackage{titlesec}
\usepackage[colorlinks=true,linkcolor=blue,citecolor=blue,urlcolor=blue]{hyperref}
\usepackage{natbib}

\titleformat{\section}{\large\bfseries}{\thesection}{1em}{}
\titleformat{\subsection}{\normalsize\bfseries}{\thesubsection}{1em}{}
\titleformat{\subsubsection}{\normalsize\bfseries\itshape}{}{0em}{}

\title{\textbf{How Few Replays Suffice:\\ Budget-Bounded Localization of Decisive Failures in LLM-Agent Trajectories}}
\author{}
\date{}

\newcommand{\ph}{\textit{[placeholder --- fill in after experiments are run]}}

\begin{document}
\maketitle
\vspace{-2em}

\begin{quote}
\small\textit{Assembled from the six confirmed stage outputs of an \texttt{auto-research-discussion} run (seed topic B2.2), post research-taste audit (\textcolor{ForestGreen}{\textbf{PASS}}). Status: \textbf{specification / paper draft --- no experiments have been run.} All numbers below are pre-registered design targets, not results.}
\end{quote}

\begin{abstract}
Diagnosing why a long LLM-agent trajectory failed increasingly relies on counterfactual replay: re-executing the trajectory from an intermediate checkpoint and checking whether the task's own verifier still fails. A decisive failure corrupts agent and environment state that every later checkpoint inherits but no earlier checkpoint carries, so this replay outcome is close to a step function of the checkpoint index --- yet existing replay-based attribution systems either spend their replay budget exhaustively, one step at a time, or prune the search heuristically at a single unexamined cost point, without ever asking how few targeted replays would have sufficed. We build a checkpointed replay harness that bills every probe in position-dependent replay units, a replay-verified fault-injected trajectory corpus balanced across plan, tool, and execution layers, and two structured localization strategies --- bisection over checkpoint indices and noise-tolerant group testing --- each de-noised by an explicit vote count and each measured against the empirical rate at which the underlying step-function assumption itself breaks down. We compare both strategies against an exhaustive per-step scan, budget-matched uniform sampling, and a zero-replay log-only judge on identical trajectories at matched replay-unit budgets, reporting per-strategy cost-accuracy curves rather than a single point estimate, so that a team can read off exactly what its own replay budget buys.
\end{abstract}

\section{Introduction}

Long LLM-agent trajectories --- dozens of chained plan, tool-call, and execution steps --- are now routine in agent development, and when one of these trajectories fails, someone has to find the one step that decided the outcome. The traditional recourse is to read the interaction log by hand, but as trajectories grow longer this manual inspection becomes prohibitively expensive and inconsistent, and it settles for correlation: a suspicious-looking step is not the same as a step that was actually decisive. The causally sound alternative is counterfactual replay --- re-executing the trajectory from an intermediate checkpoint and checking whether the task's own verifier still fails --- because only re-execution can show that a hypothesized failure point was actually responsible, rather than merely nearby.

Yet every replay re-invokes real LLM and tool calls, so replay count is itself a scarce, costly resource, and the replay-based labeling systems that generate causally verified labels treat it as an afterthought rather than a rationed budget. They replay one step at a time across the entire trajectory to build a training corpus of verified failure labels, spending a number of replays linear in trajectory length without ever asking how few replays would have sufficed to answer the one question that matters at debug time: which step was decisive, and in which layer --- plan, tool, or execution --- did it occur. A decisive failure corrupts agent and environment state that every later checkpoint inherits but no earlier checkpoint carries, which means the replay outcome is close to a step function of the checkpoint index; if that structure can be exploited by a search strategy that probes only a handful of checkpoints, causally verified attribution stops being an expensive offline labeling step and becomes cheap enough to run every time a trajectory fails.

The nearest existing work approaches this gap from adjacent but distinct angles. Systems built around exhaustive replay, such as AgenTracer \citep{AgenTracer2025}, do generate causally verified failure labels --- by replacing one step at a time and checking whether the trajectory's outcome flips --- but this exhaustive scan is treated purely as an offline corpus-construction cost, with no characterization anywhere of how much of that cost a smarter search could save. Log-only attribution systems, such as Who\&When \citep{WhoWhen2025}, sidestep replay cost entirely by having an LLM judge read the transcript, but pay for it in causal validity: their own reported results show step-level attribution is substantially worse than agent-level attribution, precisely because a judge that never re-executes anything cannot verify a hypothesis, only guess at one. Intervention-driven systems such as DoVer \citep{DoVer2025} sit between these two: they replay to verify a hypothesis, but the hypothesis itself comes from an LLM's guess about where the fault lies, and no budget-optimal search over the space of checkpoints is ever run --- the question of how few replays a structured search would need never actually gets asked. Two more recent developments sharpen, rather than close, this gap. Instrument-scale fault injection has arrived: Who\&When Pro \citep{WhoWhenPro2607} constructs 12{,}326 injected-failure trajectories with golden labels by exactly replaying a successful prefix and injecting a failure at the boundary --- but it uses that corpus to benchmark log-reading attributors, never to ask how few replays a search over checkpoints would need. And a 2026 line of efficiency-motivated attribution systems --- CHIEF's oracle-guided backtracking over a causal graph \citep{CHIEF2602}, FALAT's dependency-guided search \citep{FALAT2606}, HPFA's hypergraph comparison against a successful reference path \citep{HPFA2608}, MARS's budget-trimmed tree search over repairs \citep{MARS2607} --- now treats attribution or repair cost as a first-class concern, yet each reports the single cost operating point its own pipeline lands on; none yields the thing a provisioning decision needs, a cost-accuracy curve across strategies at matched replay spend.

Turning that question into a measurable one surfaces its own difficulties. Because a single replay re-executes an entire trajectory suffix, its cost is not a fixed per-probe unit but scales with how much of the trajectory remains --- so any claim about savings has to be billed in a cost unit that reflects this, not in raw probe counts that would misprice strategies that probe at different positions. The very step-function structure a search strategy would exploit is only approximately true: if an intermediate step can partially recover from an earlier fault, the replay signal stops being a clean threshold, and a search built to exploit that threshold can converge on the wrong index without any warning. Each individual replay is also a stochastic draw --- re-executing the same suffix twice under nonzero sampling temperature can yield different pass/fail verdicts --- so a localization strategy has to decide how many repeated replays a single probe needs before its answer is trustworthy, and that de-noising itself competes for the same limited budget. And because a trajectory spans three heterogeneous layers with different step semantics, a decisive failure can manifest across a layer boundary --- a bad early plan decision surfacing only in a later tool call --- which complicates how any search strategy should partition the space of checkpoints it searches over.

Our approach turns each of these obstacles into a specific measurement design rather than a hedge. Every replay is billed in position-dependent replay units --- resumed LLM calls, not probe counts --- so a savings claim reflects what a replay actually costs, wherever in the trajectory it is taken. A bisection search exploits the checkpoint-index step function directly, while logging the full empirical replay profile so that the rate at which that step-function assumption breaks down is measured, not assumed away, and can be checked as the specific predictor of where the bisection shortcut still holds. Each probe is de-noised by an explicit, swept number of repeated replays rather than a hidden constant, turning oracle noise from an assumption into a measured quantity. A noise-tolerant group-testing search complements bisection by partitioning the checkpoint space with built-in redundancy, including partitions that cross layer boundaries, so that attribution is not confined to contiguous runs of same-layer steps. And a replay-verified, fault-injected trajectory corpus, balanced across layers, lets every strategy --- bisection, group testing, exhaustive replay, uniform sampling, and a zero-replay log-only judge --- be compared on identical trajectories at identical replay-unit budgets, which is the shared comparison no cited system's own evaluation currently provides.

We evaluate this design on trajectories with programmatically injected, replay-verified failure locations spanning plan, tool, and execution layers in a deterministic household-task simulator and a cached-web tool-use domain, plus a human-annotated natural-failure slice, comparing every strategy's cost-accuracy curve across a shared budget grid rather than reporting a single accuracy number at a single budget.

In summary, we contribute:
\begin{itemize}[leftmargin=1.4em]
\item An identification and direct accounting of the position-dependent cost of counterfactual replay, in resumed-LLM-call units rather than raw probe counts, showing that a cost-accuracy comparison which ignores where in a trajectory a probe lands risks mispricing every localization strategy it compares.
\item A measured, falsifiable turn on the assumption that makes bisection efficient --- the monotonicity-violation rate of the replay signal --- tested directly as the predictor of where the bisection shortcut holds and where it silently breaks, rather than left as an unstated risk.
\item A characterization of the noise regime in which a redundancy-based group-testing search overtakes a vote-multiplied bisection search, and the layer regime in which crossing plan/tool/execution boundaries in a group-testing partition helps localize decisive failures that contiguous partitions miss.
\item The first head-to-head, matched-budget cost-accuracy comparison of structured replay search against exhaustive per-step replay, budget-matched uniform sampling, and zero-replay log-only judging on the same trajectories --- establishing whether, and by how much, index structure and replay access are each doing real work.
\end{itemize}

\section{Problem Definition}

\noindent\textbf{Formal setup.} Consider a trajectory $\tau = (s_1,\dots,s_n)$ that has already failed, where each step $s_i$ is one LLM call together with its tool result, and each step boundary has a checkpoint $c_i$ of agent and environment state. A replay oracle $\mathrm{replay}(k,\mathit{seed})$ re-executes steps $k,\dots,n$ from checkpoint $c_k$ --- with the recorded prefix $1,\dots,k-1$ fixed and its tool calls served from cache --- and returns the task's own verifier outcome $g(k) \in \{0,1\}$. The decisive failure step $d$ is the earliest checkpoint index $k$ at which $\mathrm{replay}(k,\cdot)$ with the original action passes while the observed or injected corruption fails --- an outcome-flipping counterfactual --- and $\ell(d) \in \{\text{plan}, \text{tool}, \text{execution}\}$ is its layer. Given a fixed replay-unit budget $B$ per trajectory, where a replay unit is one resumed LLM call and every probe (including repeated de-noising votes) is billed in this unit, the object of study is the accuracy with which a localization strategy can predict $(d,\ell(d))$ as a function of $B$ --- reported as a cost-accuracy curve, not a single point estimate.

\noindent\textbf{Why this is hard.} Because $\mathrm{replay}(k,\cdot)$ re-executes the entire suffix from $k$ onward, its cost in resumed LLM calls scales with $n-k+1$ rather than being a fixed unit --- so any claim that a strategy ``saves replays'' is only meaningful once the accounting is stated in these position-dependent units rather than in raw probe counts, which would misprice strategies that probe at different positions in the trajectory. The predicate $g(k)$ is only \emph{approximately} a step function of $k$: it is exactly the state-inheritance property --- that a decisive failure corrupts state every later checkpoint inherits and no earlier checkpoint carries --- that gives $g(k)$ its near-monotone structure, but when an intermediate step has a nonzero probability of recovering from an earlier fault, $g(k)$ can violate that monotonicity, and a search strategy built on the monotone assumption can converge on the wrong index with no signal that it has done so. $g(k)$ is also a stochastic Bernoulli read rather than a deterministic one: replaying the same suffix twice under nonzero sampling temperature can disagree, so any single probe's answer must be de-noised by repeated replay, and the number of repetitions spent de-noising competes for the same fixed budget $B$ as the search itself. Finally, because $d$ and $\ell(d)$ range over three heterogeneous layers with different step semantics and checkpoint granularity, a search strategy's partition of the index space --- contiguous runs versus partitions that cross layer boundaries --- has to be chosen without knowing in advance whether the decisive failure's cause and its observable symptom sit in the same layer.

\section{Method Overview}

At a conceptual level, the study is a single measurement pipeline with five cooperating parts, built around one shared primitive: the checkpointed suffix-replay predicate $g(k)$, billed honestly in position-dependent replay units. Bisection exploits $g(k)$'s near-monotone structure directly, while a companion diagnostic measures the rate at which that structure actually holds. A vote-based de-noising layer turns oracle noise from an assumption into a swept experimental factor. A group-testing search extends localization across layer boundaries using redundancy rather than a monotonicity assumption. And a replay-verified, layer-balanced fault-injected corpus lets every strategy --- the two structured searches, an exhaustive per-step scan, uniform sampling, and a zero-replay judge --- be compared on identical trajectories at identical budgets.

The \textbf{position-dependent replay-unit accounting} module bills every probe, including repeated votes, in resumed-LLM-call units rather than treating one probe as a fixed-cost unit, and logs probe count and measured dollar cost in parallel as secondary axes. The \textbf{bisection localizer and monotonicity diagnostic} binary-searches checkpoint indices using $g(k)$, while logging the full empirical $g$-profile so the monotonicity-violation rate can be computed and checked as the predictor of where the search shortcut holds. \textbf{Vote-based de-noising} stabilizes each probe by majority vote over an explicit, swept number of repeated replays rather than a hidden constant. \textbf{Noise-tolerant, layer-crossing group testing} partitions the checkpoint index space --- including partitions that cross plan/tool/execution boundaries --- and uses built-in redundancy to tolerate a noisy oracle where bisection alone cannot. The \textbf{matched-budget comparison corpus and baseline suite} provides replay-verified, layer-balanced ground truth and runs every strategy on the same trajectories at the same budget grid, the shared comparison this paper's gap requires.

\section{Methodology}

\subsection{Position-Dependent Replay-Unit Accounting}
\textbf{Challenge.} A single replay's cost is not a fixed per-probe unit but scales with how much of the trajectory suffix must be re-executed, so a savings claim that ignores where in the trajectory a probe lands risks mispricing every strategy it compares.

\textbf{Insight.} A cost metric billed in resumed LLM calls, rather than raw probe count, makes ``savings'' a claim about compute actually spent, and lets that claim be checked for stability against probe count and measured dollar cost as independent accounting axes.

\textbf{Method.} Every invocation of the replay oracle, including repeated de-noising votes, is logged and billed as $(n-k+1)$ resumed LLM calls, where $k$ is the checkpoint replayed from and $n$ is the trajectory length. This is the scheduled cost: a re-sampled suffix need not take exactly $n-k+1$ steps to terminate, and the measured dollar-cost axis captures any deviation between scheduled and realized spend. A trajectory's total replay cost for a strategy is the sum of these unit costs across every probe it issues. Probe count and a measured API dollar cost are logged in parallel for every trajectory, and a strategy ranking is reported as robust only if it is stable across all three axes.

\subsection{Bisection Localizer and the Monotonicity Diagnostic}
\textbf{Challenge.} Bisection's entire efficiency advantage depends on $g(k)$ being a clean step function of the checkpoint index, but intermediate steps can partially recover from an earlier fault, silently breaking that assumption with no built-in warning.

\textbf{Insight.} The monotonicity assumption cannot be guaranteed analytically, but its violation rate can be measured directly from the same probes bisection already issues, turning a silent failure mode into a reported, falsifiable diagnostic.

\textbf{Method.} Binary search over checkpoint indices uses the suffix-replay predicate $g(k)$ to narrow the candidate range for the decisive step in each iteration, converging in a number of probes logarithmic in trajectory length. Every probed $g(k)$ value is retained in a per-trajectory profile; the monotonicity-violation rate is computed as the fraction of trajectories whose empirical profile contains any pair $k<k'$ with $g(k)=0$ and $g(k')=1$, and this rate is reported alongside bisection's accuracy as its load-bearing predictor.

\subsection{Vote-Based De-Noising}
\textbf{Challenge.} A single replay is a stochastic read of $g(k)$ --- re-executing the identical suffix twice under nonzero sampling temperature can disagree --- so any one probe's pass/fail answer cannot be trusted without repetition, but repetition itself consumes budget.

\textbf{Insight.} Treating the number of repeated replays per probe as an explicit, swept experimental factor rather than a fixed assumption turns ``is the oracle noisy enough to matter'' from a hidden constant into a measured, reportable quantity.

\textbf{Method.} Each probe's $g(k)$ read is taken as the majority vote over $r$ repeated replays at that checkpoint, with $r$ swept over $\{1,3,5\}$ as an independent experimental factor. The single-replay flip rate --- the probability that two independent replays of the same checkpoint disagree --- is estimated per trajectory and reported alongside the vote count needed to stabilize the resulting probe.

\subsection{Noise-Tolerant, Layer-Crossing Group Testing}
\textbf{Challenge.} A search confined to contiguous runs of checkpoints cannot attribute a failure whose cause and observable symptom sit in different layers, and an oracle read that is both noisy and position-dependent needs a de-noising mechanism beyond simply repeating individual probes.

\textbf{Insight.} Redundancy across groups of checkpoints --- rather than repetition of a single checkpoint --- is the standard combinatorial mechanism for tolerating a noisy oracle, and defining those groups over the whole index space, not just contiguous spans, is what lets a search strategy attribute across a layer boundary.

\textbf{Method.} The step-index space is partitioned into groups, either contiguous or explicitly layer-crossing, and each group is tested via a multi-step intervention (replacing the group's steps with recorded-correct actions where available, or re-sampling otherwise) with redundancy built into the group design to tolerate a noisy $g(k)$ read; the implicated group is then refined recursively. Both a fixed-noise, varied-partition sweep and a fixed-partition, varied-noise sweep are run from this single design, so that the noise-crossover claim and the layer-crossing-design claim are each tested from the same underlying mechanism.

\subsection{Matched-Budget Comparison Corpus and Baseline Suite}
\textbf{Challenge.} No existing replay-based attribution system's own evaluation compares a structured search strategy, an exhaustive scan, uniform sampling, and a zero-replay judge on the same trajectories at the same replay-unit budgets, so no accuracy-vs-replay-budget characterization currently exists to compare against.

\textbf{Insight.} A characterization is only informative if every strategy is run against ground truth that is verified independently of any one strategy, on identical trajectories, at identical budgets --- otherwise a ``savings'' claim could be an artifact of which trajectories or which budgets happened to be compared.

\textbf{Method.} Following programmatic fault-injection practice, one step of a successful trajectory is perturbed (a subgoal in the plan layer, an argument in the tool layer, an observation in the execution layer), the perturbed run is replayed downstream, and the trajectory is retained only if the perturbed run fails and a replay from the injected step with the original action succeeds --- yielding a replay-verified decisive step index and layer per trajectory, balanced across the three layers. Every localization strategy --- bisection, group testing, uniform-random sampling, an exhaustive per-step scan, and a zero-replay log-only judge --- is then run on the identical corpus at the identical replay-unit budget grid.

\section{Experiments}

\subsection{Experimental Setup}

\textbf{Datasets.} ALFWorld (deterministic household-task simulator) --- used for the pilot stage because its full determinism lets replay reproducibility be gated at $\geq$95\% before any strategy comparison is trusted. GAIA cached-web subset (tool-use tasks, snapshotted web content for deterministic replay) --- introduces tool-layer failures with real, though cached, tool-call variability absent from ALFWorld's synthetic simulator, letting the layer-crossing-design and structure-vs-uniform-and-judge research questions be tested under a genuine tool-use environment rather than only a household-task simulator. Who\&When \citep{WhoWhen2025} human-annotated failure logs --- the natural-failure evaluation slice, re-hosted in the replay harness where re-executable, exposing any gap between injected-fault and natural-failure localization accuracy.

\textbf{Baselines} (sanity-check and isolation comparators, not competitors to beat --- this is a measurement study, and no research question here is framed as ``our procedure beats baseline X''). Budget-matched uniform-random checkpoint sampling isolates whether any advantage comes from the index structure bisection and group testing exploit, or merely from having replay access and spending a budget at all. An exhaustive per-step counterfactual scan, following current fault-injection-corpus practice \citep{AgenTracer2025}, serves as the accuracy ceiling, not a rival to beat. A zero-replay log-only LLM judge, in Who\&When's \citep{WhoWhen2025} all-at-once and step-by-step configurations, is the no-mechanism control: if it matches replay-based accuracy, the entire replay enterprise is unjustified. (Who\&When's third configuration, binary search, bisects the log context shown to the judge rather than any execution --- a context-management variant of the same log-only mechanism; the two canonical configurations are reported here.) A DoVer-style \citep{DoVer2025} hypothesis-then-verify baseline checks whether LLM-proposed hypotheses plus replay verification match a budget-blind structured search, without adopting DoVer's own reported numbers. Separately from these comparators, a null-by-construction placebo (B0b, \texttt{paper.md}'s mandatory run-validity gate rather than a research-question arm) is scored before any hypothesis: a label-shuffle placebo scores every strategy's predictions against ground-truth indices permuted across trajectories, which must land inside the analytic chance band, and a sham-injection arm applies the fault-injection machinery at a step whose corruption verifiably does not flip the outcome, on which replay-based strategies must read ``no fault'' and the judge's predicted index must be at chance with respect to the sham site --- a non-null B0b means the pipeline measures an artifact and the whole run is invalid.

\textbf{Metrics.} Exact-step localization accuracy and step accuracy at a $\pm$2-step tolerance; layer accuracy; replay cost in replay units (resumed LLM calls, position-weighted), with probe count and measured API dollar cost as secondary axes; the monotonicity-violation rate of the empirical $g(k)$ profile; and normalized area under the accuracy-vs-log-budget curve over the shared budget grid.

\textbf{Research questions:}
\begin{enumerate}[leftmargin=1.6em]
\item RQ1 --- Does billing replay cost in resumed-LLM-call units reveal that structured strategies' savings over exhaustive scanning depend on where the decisive step falls?
\item RQ2 --- Does bisection reach exhaustive-comparable localization accuracy at a fraction of its replay-unit cost, on the pooled corpus?
\item RQ3 --- Does the monotonicity-violation rate predict where bisection's accuracy advantage over uniform sampling shrinks?
\item RQ4 --- How many repeated replays per probe are needed to stabilize bisection's localization accuracy?
\item RQ5 --- Is there an oracle-noise level above which group testing's accuracy-per-replay-unit exceeds vote-multiplied bisection's?
\item RQ6 --- Do layer-crossing group-testing partitions outperform contiguous partitions specifically on plan-layer decisive failures, and does the structured advantage hold in every layer?
\item RQ7 --- Do both structured strategies beat budget-matched uniform sampling and a zero-replay judge at a shared budget?
\end{enumerate}

\subsubsection{RQ1: Does billing replay cost in resumed-LLM-call units reveal that structured strategies' savings over exhaustive scanning depend on where the decisive step falls?}
\textbf{Registration status.} Per \texttt{paper.md} Step 6's registration of record, this is the pre-registered \emph{descriptive} cost-accounting analysis (A7), not a confirmatory mechanism-tier hypothesis --- it is reported alongside H1--H4 but carries no PASS/FAIL verdict of its own; the falsification condition below states what would make the position-dependent accounting claim non-load-bearing, not a hypothesis test outcome.

\textbf{Setup.} For every trajectory in the pilot and scale-up corpora, compute each strategy's total replay-unit cost (resumed LLM calls) and its exact-step accuracy, then stratify the ratio of structured-strategy cost to exhaustive-scan cost by the quartile of the trajectory in which the decisive step $d$ falls (early/mid/late).

\textbf{Expected outcome.} The replay-unit cost ratio (structured strategy / exhaustive scan) is smallest --- i.e., savings are largest --- for late-decisive-step trajectories, where a suffix replay is cheapest to begin with; savings are smaller, though still present, for early-decisive-step trajectories.

\textbf{Falsification condition.} Refuted if the cost ratio is uniform across where the decisive step falls (no quartile effect), or if a probe-count-based accounting would rank strategies differently from a replay-unit accounting --- either would mean the position-dependent cost claim is not load-bearing and a simpler probe-count metric would have sufficed.

\textbf{Hypothesis (plain language).} How much a structured search actually saves depends on where in the trajectory the decisive step sits, and only a cost metric that reflects real replay expense --- not raw probe counts --- reveals that dependence.

\textbf{Metrics.} Replay-unit cost ratio by decisive-step-position quartile; rank stability of strategies across replay-unit, probe-count, and dollar-cost accounting.

\textbf{Result.} \ph

\subsubsection{RQ2: Does bisection reach exhaustive-comparable localization accuracy at a fraction of its replay-unit cost?}
\textbf{Setup.} Pilot stage: 2 backbones $\times$ 2 conditions (bisection, exhaustive scan) $\times$ ALFWorld, 60 trajectories (20 per layer). Scale-up: $\geq$4 backbones $\times$ 2 domains (ALFWorld, GAIA cached-web subset) $\times$ 300 trajectories. \emph{Registration note:} the pre-registered criterion (\texttt{paper.md} Step 6, H1 --- the sole binding registration) is scored on the \textbf{pooled} scale-up corpus, not restricted to the monotone-profile subset; an earlier draft of this research question scoped the criterion to trajectories with a monotone empirical $g(k)$ profile, but \texttt{paper.md}'s registration explicitly supersedes that restriction, so it is dropped here. Paired comparison of exact-step accuracy (McNemar's test) and replay-unit cost (Wilcoxon signed-rank) between bisection and the exhaustive scan, $\alpha = 0.05$, Holm-corrected across backbones and domains. Cost accounting holds the de-noising vote count fixed across arms --- the exhaustive scan is billed at the same per-probe vote count $r$ as bisection (chosen per RQ4) --- so the cost ratio isolates search structure rather than a de-noising asymmetry between arms.

\textbf{Expected outcome.} On the pooled scale-up corpus, bisection reaches at least 90\% of the exhaustive scan's exact-step accuracy while consuming at most 25\% of its replay-unit cost.

\textbf{Aggregation.} PASS requires the pooled effect \emph{and} the same-direction effect (bisection within the stated accuracy/cost floor of exhaustive) in at least 3 of the 4 environment$\times$LLM-family strata, so a pooled PASS cannot be carried by one favorable stratum while the others disagree; if an entire environment stratum is lost to the harness reproduction-rate gate, the rule degrades to same-direction in 2 of 2 remaining strata and the claim's envelope narrows to the surviving environment.

\textbf{Falsification condition.} Refuted if bisection needs a replay-unit budget at or above 50\% of the exhaustive scan's to come within 5 accuracy points of it, on the pooled corpus --- that would mean the step-function structure buys far less than the logarithmic saving the design is built to exploit.

\textbf{Hypothesis (plain language).} Where the replay signal really is close to a step function of the checkpoint index, bisection captures nearly all of the exhaustive scan's accuracy for a small fraction of its cost.

\textbf{Metrics.} Exact-step accuracy and 95\% CI (pooled corpus); replay-unit cost ratio (bisection / exhaustive).

\textbf{Result.} \ph

\subsubsection{RQ3: Does the monotonicity-violation rate predict where bisection's accuracy advantage over uniform sampling shrinks?}
\textbf{Registration status.} Per \texttt{paper.md} Step 6's registration of record, this is the pre-registered \emph{descriptive} mechanism analysis under A3/F1, not a confirmatory mechanism-tier hypothesis --- it is reported as a stratified diagnostic that informs how H1's fallback (F1) is worded if H1 fails, not as an independently PASS/FAIL-scored claim.

\textbf{Setup.} Compute the monotonicity-violation rate per trajectory from bisection's logged $g(k)$ profile; stratify trajectories into violating vs. non-violating strata; compare bisection's accuracy advantage over budget-matched uniform sampling within each stratum via paired McNemar, Holm-corrected across strata.

\textbf{Expected outcome.} In the monotonicity-violating stratum, bisection's accuracy advantage over uniform sampling shrinks to within 5 percentage points; in the non-violating stratum, the advantage remains large (consistent with RQ2's expected margin).

\textbf{Falsification condition.} Refuted if the advantage is statistically indistinguishable between the violating and non-violating strata --- that would mean the monotonicity-violation rate carries no predictive information about where bisection's shortcut holds, contradicting the diagnostic's intended purpose.

\textbf{Hypothesis (plain language).} The monotonicity-violation rate is not just a reported statistic but the specific, checkable reason bisection sometimes underperforms --- not an unrelated source of noise.

\textbf{Metrics.} Per-stratum accuracy advantage (bisection $-$ uniform) and 95\% CI; monotonicity-violation rate distribution.

\textbf{Result.} \ph

\subsubsection{RQ4: How many repeated replays per probe are needed to stabilize bisection's localization accuracy?}
\textbf{Setup.} Run bisection with $r \in \{1,3,5\}$ repeated replays per probe at sampling temperature $\approx 0.7$ on the same trajectories; compare exact-step accuracy across $r$ via paired Wilcoxon signed-rank tests; estimate the single-replay flip rate per trajectory. This research question carries conjunct (i) of H4 (\texttt{paper.md} Step 6): H4's $\alpha=0.05$ is Holm-split across its two conjuncts (per-conjunct $\alpha=0.025$), with a further Holm correction across the three pairwise $r$-comparisons within this conjunct.

\textbf{Expected outcome.} $r=3$ stabilizes bisection's accuracy to within 3 percentage points of $r=5$, indicating diminishing returns to further de-noising beyond $r=3$.

\textbf{Falsification condition.} Refuted if accuracy at $r=1$ is statistically indistinguishable from $r=5$ (the oracle is effectively deterministic and de-noising was unnecessary --- a scoping result, not a failure), or if the gap between $r=3$ and $r=5$ exceeds 3 points (more votes are required than assumed, raising the effective cost of bisection beyond the RQ2 estimate).

\textbf{Hypothesis (plain language).} A modest, fixed number of repeated replays per probe is enough to make bisection's noisy oracle trustworthy, without the vote count itself eating the budget advantage.

\textbf{Metrics.} Exact-step accuracy vs. $r$; single-replay flip rate per trajectory.

\textbf{Result.} \ph

\subsubsection{RQ5: Is there an oracle-noise level above which group testing's accuracy-per-replay-unit exceeds vote-multiplied bisection's?}
\textbf{Setup.} At a fixed group-testing partition design (contiguous), sweep the single-replay flip rate (via sampling temperature) and compare accuracy-per-replay-unit between group testing and vote-multiplied bisection ($r$ chosen per RQ4) across the sweep; locate the crossover point with a 95\% bootstrap confidence interval (1{,}000 resamples) on the flip-rate axis, required to exclude ``no crossover within the observed flip-rate range.'' This research question carries conjunct (ii) of H4 (\texttt{paper.md} Step 6), scored at the per-conjunct $\alpha=0.025$ (Holm split of H4's $\alpha=0.05$ across its two conjuncts).

\textbf{Expected outcome.} A flip rate exists (expected in the neighborhood of 15\%) above which group testing's accuracy-per-replay-unit exceeds bisection's, because redundancy across groups amortizes noise-tolerance cost more efficiently than repeated single-checkpoint votes at high noise.

\textbf{$p$-axis support pre-condition.} This conjunct is scorable only if the observed per-trajectory flip rates actually span the 15\% threshold region under the swept sampling temperatures. Both evaluation environments (ALFWorld, GAIA cached-web) were chosen specifically for deterministic, reproducible replay (the $\geq$95\% reproduction gate), so this is a live risk rather than a formality: if no trajectories in the corpus reach the threshold flip rate at any swept temperature, this conjunct is pre-declared \textbf{INCONCLUSIVE (insufficient noise support)}, not FAIL, and is reported as untested in this noise regime rather than as a retracted claim.

\textbf{Falsification condition.} Refuted if flip rates do span the threshold region but no flip rate in that range shows group testing exceeding bisection's accuracy-per-replay-unit --- that would mean vote-multiplied bisection dominates across the whole tested, adequately-supported noise range and the group-testing noise-tolerance claim does not hold.

\textbf{Hypothesis (plain language).} Past a certain noise level, spreading redundancy across a group of checkpoints is a more efficient way to tolerate oracle noise than repeating votes at a single checkpoint.

\textbf{Metrics.} Accuracy-per-replay-unit vs. flip rate for both strategies; bootstrap CI on the crossover flip rate.

\textbf{Result.} \ph

\subsubsection{RQ6: Do layer-crossing group-testing partitions outperform contiguous partitions specifically on plan-layer decisive failures, and does the structured advantage over uniform sampling hold in every layer?}
\textbf{Setup.} At a fixed noise level, compare contiguous vs. explicitly layer-crossing group-testing partition designs on exact-step and layer accuracy, stratified by the decisive step's true layer; paired comparison per layer stratum, Holm-corrected across the three layers ($\alpha=0.017$ per layer after correction). Balanced corpus construction guarantees $\approx$90--100 trajectories per layer at $n=300$. This research question carries \textbf{both conjuncts of H3} (\texttt{paper.md} Step 6): (i) the $\geq$10-point structured-vs-uniform advantage holding per layer, and (ii) the layer-crossing-vs-contiguous margin on plan-layer faults specifically.

\textbf{Expected outcome.} Layer-crossing partitions outperform contiguous partitions by at least 5 percentage points on plan-layer decisive failures specifically, with little or no difference on tool- or execution-layer failures; independently, the $\geq$10-point structured-over-uniform advantage (RQ7's pooled effect) holds within each layer stratum.

\textbf{Aggregation.} At the stated per-layer power ($\approx$65--70\% at $\alpha=0.017$, below the pooled check's $>$80\%), requiring all 3 layers to independently clear significance would mostly measure power rather than truth ($\approx$0.65$^3$--0.70$^3\approx$0.27--0.34 even if the true effect holds everywhere). H3's first conjunct therefore PASSes iff the $\geq$10-point advantage holds in at least 2 of the adequately-powered layers \emph{and} no adequately-powered layer shows a statistically significant reversal; under-powered layers are pre-declared INCONCLUSIVE for that layer and excluded from the denominator, and if fewer than 2 layers are adequately powered the first conjunct as a whole is INCONCLUSIVE. The second conjunct (layer-crossing margin on plan-layer faults) is scored separately and unchanged by this rule; H3 PASSes only if both conjuncts PASS.

\textbf{Falsification condition.} The second conjunct is refuted if layer-crossing partitions show no advantage on plan-layer failures, or show an equal or larger advantage on tool- or execution-layer failures --- either would mean the layer-crossing design's benefit is not specifically tied to plan-layer failures manifesting downstream, contradicting the stated mechanism. The first conjunct is refuted (not merely INCONCLUSIVE) only if an adequately-powered layer shows a statistically significant reversal (uniform outperforming structured).

\textbf{Hypothesis (plain language).} Because a bad plan decision is the failure type most likely to surface only in a later tool or execution step, only plan-layer failures specifically benefit from a search partition that is allowed to cross layer boundaries.

\textbf{Metrics.} Per-layer exact-step and layer accuracy, contiguous vs. layer-crossing partitions.

\textbf{Result.} \ph

\subsubsection{RQ7: Do both structured strategies beat budget-matched uniform sampling and a zero-replay judge at a shared budget?}
\textbf{Setup.} At a shared confirmatory replay-unit budget of one-quarter of the exhaustive scan's per-trajectory replay-unit cost --- $B = C_{\mathrm{exh}}/4$, where $C_{\mathrm{exh}} = n(n+1)/2$ resumed calls for a length-$n$ trajectory --- compare bisection and group testing against budget-matched uniform sampling and the zero-replay log-only judge on exact-step accuracy, paired McNemar per comparison, Holm-corrected across the four pairwise comparisons; $n=300$ trajectories in the scale-up corpus, sized to give $>$80\% power to detect a 10-percentage-point paired accuracy difference at $\alpha=0.05$ assuming $\approx$20\% discordance.

\textbf{Registration amendment (2026-08-15).} The earlier wording of this budget point --- ``a replay budget of one-quarter the trajectory length,'' $B = n/4$ replay units, with a descriptive grid $B \in \{2,\dots,64,\text{exhaustive}\}$ replay units --- is incoherent under the position-dependent accounting this paper itself registers: at the corpus's mean length $n \approx 30$, $B \approx 8$ resumed calls cannot fund a single mid-trajectory probe (cost $n/2+1 \approx 16$), and bisection's own full-run expected cost ($\approx (n/2)\log_2 n \approx 75$ units at $r=1$) exceeds the raw grid's maximum of 64 --- the raw-unit wording silently re-imports the probe-count mispricing RQ1 exists to expose. The budget axis is therefore restated in normalized units: fractions of the exhaustive cost, with descriptive grid $B/C_{\mathrm{exh}} \in \{1/64, 1/32, 1/16, 1/8, 1/4, 1/2, 1\}$ and the confirmatory point at $B/C_{\mathrm{exh}} = 1/4$. This preserves the registered intent --- $n/4$ uniformly placed probes cost $\approx (n/4)\cdot(n/2) \approx C_{\mathrm{exh}}/4$ --- while removing the unit ambiguity; \texttt{paper.md} Step 6 (H2) requires the matching amendment before any Phase B run (logged in \texttt{pipeline.md}).

\textbf{Expected outcome.} Both bisection and group testing exceed budget-matched uniform sampling's exact-step accuracy by at least 15 percentage points, and the zero-replay judge sits at least 20 percentage points below the best structured strategy.

\textbf{Aggregation.} As with RQ2/H1, PASS is scored on the pooled n=300 corpus at the single pre-registered budget $B=C_{\mathrm{exh}}/4$ and additionally requires the same-sign effect in at least 3 of the 4 environment$\times$LLM-family strata; the full budget-grid curve ($B/C_{\mathrm{exh}}\in\{1/64,\dots,1/2,1\}$) is reported descriptively at every grid point, but only the $B=C_{\mathrm{exh}}/4$ point is confirmatory.

\textbf{Falsification condition.} Refuted (structure adds nothing) if either structured strategy fails to beat budget-matched uniform sampling by a statistically significant margin; refuted (replay is pointless) if the zero-replay judge's accuracy is statistically indistinguishable from the best structured strategy's.

\textbf{Hypothesis (plain language).} At a shared, realistic budget, both the index structure a search strategy exploits and the fact of having replay access at all are each doing real, separately necessary work --- neither alone would explain the full advantage.

\textbf{Metrics.} Exact-step accuracy per strategy at $B=C_{\mathrm{exh}}/4$, pairwise accuracy gaps and 95\% CIs.

\textbf{Result.} \ph

\vspace{0.5em}
\noindent\textbf{Further evaluation.} With more time, a natural extension is to test whether the cost-accuracy curves transfer to multi-agent orchestration frameworks or to trajectories an order of magnitude longer than those sampled here --- explicitly out of scope in this design because both would require re-estimating the position-dependent cost model and the monotonicity-violation rate on a different trajectory-length distribution, rather than being a direct extension of the budget grid used above.

\section{Related Work}

The closest prior work to this paper is AgenTracer \citep{AgenTracer2025} and DoVer \citep{DoVer2025}, and it is worth stating the delta plainly before surveying the wider landscape: AgenTracer already generates causally verified failure labels through counterfactual replay, and DoVer already replays from a hypothesized intervention point to verify it --- but neither asks, nor answers, how few replays a structured search over checkpoint indices would need to reach the same accuracy, which is precisely the budget-bounded characterization this paper is built to produce. One concession must be made up front: Who\&When Pro \citep{WhoWhenPro2607} has since scaled replay-verified fault injection to 12{,}326 golden-label trajectories across 26 benchmarks, so the injected-fault corpus built here is no longer a first-at-scale instrument --- the corpus delta this paper retains, and the reason the estimand still requires its own corpus, are stated precisely below.

\textbf{Replay-based failure attribution and labeling.} AgenTracer \citep{AgenTracer2025} generates causally grounded decisive-error labels by replaying a failed trajectory one step at a time and checking whether the outcome flips, training a specialized tracer on the resulting corpus; DoVer \citep{DoVer2025} treats attribution as a testable causal hypothesis, replaying from an LLM-suggested intervention point to check whether the outcome recovers; Who\&When \citep{WhoWhen2025} instead reads the transcript with an LLM judge, in all-at-once, step-by-step, and binary-search configurations, without replaying anything --- its binary-search variant bisects the log context handed to the judge, not an execution, so it shares the judge family's causal blind spot rather than anticipating the replay bisection studied here. All three establish that a failed trajectory's decisive step can in principle be identified, but AgenTracer's labeling replay is exhaustive by construction, DoVer's replay verifies only whichever hypothesis an LLM proposes rather than searching the checkpoint space itself, and Who\&When's log-only judge never re-executes anything and --- by its own reported results --- is substantially worse at step-level than agent-level attribution. None characterizes an accuracy-vs-replay-budget tradeoff, the gap this paper's cost-accounted search strategies are built to close. A closely related but distinct line of work formulates single-trajectory attribution as a causal-intervention estimation problem rather than a search-strategy problem: Causal Agent Replay \citep{CausalAgentReplay2606} models a trajectory as a structural causal model and estimates per-step causal magnitude (a contrastive point-of-commitment locus, and a budget-bounded Monte-Carlo Shapley estimator for multi-step interactions), validated against synthetic ground truth. Its ``budget-bounded'' qualifier governs the internal Monte-Carlo sample count of its Shapley estimator, not a comparison of localization search strategies at matched replay-unit budgets, so it is adjacent rather than a method-twin: it asks which estimator best attributes a fixed trajectory's causal structure, not how few replays a structured search needs to find the decisive step in the first place. Three further replay-adjacent systems complete this line. Who\&When Pro \citep{WhoWhenPro2607} scales the injection route to instrument grade --- 12{,}326 injected-failure trajectories with golden labels, each built by exactly replaying a successful prefix and injecting a failure at the boundary --- but its verification stops at the prefix: there is no round-trip check that the original action at the injected step still rescues the outcome, no sham-injection control, no layer balancing, and no per-step checkpoint packaging for suffix replay, which are exactly the properties the budget-bounded search estimand requires of its corpus; and its evaluation benchmarks log-reading attributors at a fixed, budget-free operating point. FAMAS \citep{FAMAS2509} attributes by spectrum analysis over repeated full executions --- replay-based, but with the number of executions never priced or varied as a budget. AgentDebug \citep{AgentDebug2509} isolates root-cause steps with a counterfactual test and contributes AgentErrorBench, annotated failure trajectories on ALFWorld, GAIA, and WebShop --- the same environments used here --- but characterizes no cost-accuracy tradeoff for the isolation itself.

\textbf{Efficiency-motivated structural search over failure traces.} A rapidly growing 2026 line makes the cost of attribution itself the motivation: CHIEF \citep{CHIEF2602} converts flat logs into a hierarchical causal graph and prunes candidate causes by oracle-guided backtracking before counterfactual screening; FALAT \citep{FALAT2606} frames attribution as dependency-guided search, checking whether correcting a candidate step recovers the expected outcome; StepFinder \citep{StepFinder2606} distills the judge into a lightweight temporal model precisely to avoid repeated expensive LLM passes; HPFA \citep{HPFA2608} names exhaustive step-by-step counterfactual testing as its motivating enemy and prunes it by comparing hyperedges against a successful reference path; MARS \citep{MARS2607} couples attribution to repair via Monte-Carlo tree search with partial rollouts to cut token spend, and releases StateMAS, 1{,}310 replayable failure trajectories. Each of these buys efficiency through structure, and each reports the single cost operating point its own pipeline lands on; none compares localization strategies on identical trajectories at matched, position-priced replay budgets, and none measures the monotone-signal assumption a budget-bounded index search would lean on. MP-Bench \citep{MPBench2603} is the construct-level caveat over this whole cluster, the present paper included: real multi-agent failures often admit multiple defensible attributions, which is why the headline numbers here are licensed only on a replay-verified single-fault corpus and the natural-failure slice is reported separately.

\textbf{Step-level credit assignment and compute allocation during generation or training.} InT \citep{InT2025} localizes the first erroneous step of a reasoning rollout via reference-grounded self-verification and fine-tunes on the resulting prefix-plus-intervention pairs; TRIM \citep{TRIM2025} allocates a strong-model token budget to the few reasoning steps a risk signal flags as likely to cascade into error; AttnRL \citep{AttnRL2025} scores each step's forward context influence to place Monte Carlo rollouts only at the most influential steps; iStar \citep{iStar2025} derives dense per-action credit from trajectory preference pairs without any replay. Each of these intervenes \emph{during} generation, training, or rollout construction --- none provides post-hoc localization of a decisive step in an already-failed, already-logged trajectory, and none reports cost-accuracy guarantees for that diagnostic task specifically, which is the setting this paper's replay harness and evaluation are built for.

\textbf{Fixed-budget query theory and classical bisection.} Best-arm identification with a fixed sampling budget \citep{BAILD2025} analyzes adaptive query allocation over stateless arms with i.i.d.\ reward draws, the formal template for ``identify a target under a fixed number of noisy queries'' --- but it does not model a sequence of causally ordered checkpoints whose query signal has approximately monotone structure and whose per-query cost varies with checkpoint position, so its guarantees do not transfer directly to trajectory localization. Classical delta debugging \citep{ZellerHildebrandt2002} recursively bisects a failure-inducing input against a cheap, deterministic, arbitrarily repeatable test, isolating a minimal failure-inducing configuration in near-logarithmic test count --- but LLM-agent replay violates all three of those founding assumptions (each replay costs real model and tool calls, outcomes are stochastic under sampling, and environments can drift), so its minimality guarantees do not carry over unmodified. This paper does not propose bisection or group testing as new search algorithms; the contribution is characterizing how these classical primitives, and the assumptions underneath them, behave once every one of those assumptions is violated by a real LLM replay oracle.

\section{Conclusion}

This work reframes counterfactual replay's own scarcest resource --- the number of times a trajectory must be re-executed --- as a measurable characterization problem rather than an assumed-necessary cost. The central methodological wager is that a decisive failure's state-inheritance structure makes the replay signal close enough to a step function of the checkpoint index that a structured search can exploit it, but that ``close enough'' has to be measured, not assumed, because recovery-induced violations of that structure are exactly where a naive search silently fails. Whichever regime dominates for current LLM agents and environments, the finding is decision-relevant: if structured replay reaches near-exhaustive accuracy at a small fraction of the cost, causally verified attribution becomes cheap enough to run in every development loop rather than reserved for offline corpus construction; if it does not, the monotonicity-violation rate this design measures becomes the practical signal for when to fall back to exhaustive replay instead. Two limitations bound this claim from the outset and are natural next steps rather than afterthoughts: ground truth here is built from programmatically injected, replay-verified single faults, and the paired natural-failure slice may reveal that diffuse, multi-causal real failures erode the accuracy this design achieves on constructed ones; and the corpus spans two sandboxed, replay-controllable environments, leaving open whether the measured cost-accuracy curves transfer to live tool ecosystems or multi-agent orchestration frameworks where faithful, deterministic re-execution cannot be guaranteed.

\bibliographystyle{plainnat}
\bibliography{refs}

\vspace{1em}
\hrule
\vspace{1em}

\section*{Pre-handoff checklist}
\begin{itemize}[leftmargin=1.4em]
\item \textbf{Organization/navigation} --- narrative flows problem $\to$ method $\to$ experiments $\to$ related work without forward references to undefined terms.
\item \textbf{Terminology discipline} --- ``replay unit,'' ``decisive step,'' ``monotonicity-violation rate,'' ``layer-crossing partition'' each defined at first use in Problem Definition, used consistently thereafter; no internal stage IDs (\texttt{sp\_*}, \texttt{m\_*}, \texttt{rq\_*}) leaked into prose --- checked directly against the discussion pipeline log.
\item \textbf{Notation with worked example} --- the formal setup section defines $(\tau, s_i, c_i, \mathrm{replay}(k,\mathit{seed}), g(k), d, \ell(d), B)$ and the Problem Definition/Methodology sections consistently reuse this notation; no separate worked numeric example is included since no numbers exist yet (would require fabricating results).
\item \textbf{RQ framing matches outcomes} --- all 7 RQs are graded/quantitative (rates, ratios, accuracy gaps), each answered in its own subsection with an explicit falsifiable expected outcome and a distinct falsification condition, not a binary yes/no.
\item \textbf{No typesetting artifacts} --- only explicit placeholder markers in Results, per the drafting convention this is a pre-registered spec, not a report.
\item \textbf{Artifact plan} --- the checkpointed replay harness and the replay-verified fault-injected corpus are the natural release artifacts; not yet built, so not claimed as available.
\end{itemize}

\section*{Hand-off}
This is a draft, not a finished paper --- experiments have not run. Natural next steps: prose polish and venue-specific LaTeX formatting once a target venue is chosen, or revising an earlier stage if piloting surfaces a problem (e.g., if the monotonicity-violation rate in RQ3 is much higher than expected).

\end{document}
